\section{Risk Analysis}
\label{sec:risk-analysis}

This section identifies key risks that may affect the design, implementation, and evaluation of the proposed method, together with mitigation measures. Likelihood and impact are assessed qualitatively as \emph{Low}, \emph{Medium}, or \emph{High}.

\begin{table}[h!tb]
\centering
\caption{Key project risks and mitigation measures.}
\label{tab:risks}
\begin{tabular}{p{0.20\linewidth} p{0.12\linewidth} p{0.12\linewidth} p{0.48\linewidth}}
\toprule
\textbf{Risk} & \textbf{Likelihood} & \textbf{Impact} & \textbf{Mitigation} \\
\midrule

Limited exploration coverage in GUI State Models &
High & High &
Use multiple TESTAR runs and, where relevant, multiple configurations, and aggregate execution evidence across models. Maintain a clear split between models used for method development and models reserved for final evaluation. Interpret \emph{not covered} explicitly as \emph{not covered in the available model evidence}, and prioritize user stories with observable GUI outcomes. \\

Ambiguity or underspecification in user stories &
Medium & Medium &
Represent user stories in a structured form with essential steps and acceptance criteria. Define a small glossary for domain terms and GUI synonyms, and document interpretation rules and assumptions for ambiguous stories. \\

Noisy or incomplete execution evidence &
Medium & High &
Apply preprocessing steps such as label normalization, deduplication, and filtering of unstable identifiers. Prefer robust evidence signals, for example combinations of action type, stable UI attributes, and surrounding context. Include error analysis to identify common sources of mismatch. \\

Reference set creation is too costly or inconsistent &
Medium & High &
Limit the reference set to a representative subset of stories and model evidence. Use explicit labeling guidelines, record rationale for decisions, and, where feasible, have a second reviewer label a subset to assess consistency. Resolve disagreements through documented adjudication and report uncertainty where evidence is ambiguous. \\

Prototype scope or implementation effort exceeds available time &
Medium & High &
Timebox development and deliver a minimal end-to-end pipeline early, from model import and story parsing to trace generation and coverage reporting. Treat advanced strategies as optional extensions, and keep components modular to reduce rework if representations or schemas change. \\

Overfitting to a single SUT or exploration configuration &
Medium & Medium &
Evaluate the method on multiple runs or configurations where possible, avoid hard-coded SUT-specific rules, and document which design choices are tool- or SUT-specific. If time permits, perform a limited validation on a second small SUT. \\

Limited interpretability or reproducibility of tracing decisions &
Medium & Medium &
Prefer deterministic baselines where possible, version-control code, configurations, and input artifacts, and record parameters and thresholds explicitly in the repository and thesis. When multiple strategies are evaluated, include a comparison between baseline and enhanced variants. \\

\bottomrule
\end{tabular}
\end{table}

% \section{Risk Analysis}
% \label{sec:risk-analysis}

% This section identifies key risks that may affect the design, implementation, and evaluation of the proposed method and lists mitigation measures.
% Likelihood and impact are assessed qualitatively as \emph{Low/Medium/High}.

% \begin{table}[h!tb]
% \centering
% \caption{Key project risks and mitigations.}
% \label{tab:risks}
% \begin{tabular}{p{0.18\linewidth} p{0.14\linewidth} p{0.12\linewidth} p{0.48\linewidth}}
% \toprule
% \textbf{Risk} & \textbf{Likelihood} & \textbf{Impact} & \textbf{Mitigation} \\
% \midrule

% Limited exploration coverage in GUI State Models &
% High & High &
% Use multiple runs and configurations and aggregate evidence across models; keep an explicit development/evaluation split (development models for iteration, frozen models for final scoring);
% interpret ``not covered'' as ``not covered in available model evidence''; prioritize user stories with observable UI outcomes. \\

% Ambiguity or underspecification in user stories &
% Medium & Medium &
% Represent stories in a structured format with essential steps/acceptance criteria; define a small glossary for domain terms and UI synonyms; document interpretation rules and assumptions per story. \\

% Noisy or incomplete model evidence (labels, duplicates, unstable identifiers) &
% Medium & High &
% Normalize labels and tokens; apply deduplication and filtering heuristics; favor robust evidence signals (e.g., action type + stable attributes + surrounding context);
% include an error analysis that reports typical mismatch causes. \\

% Reference set creation is too costly or inconsistent &
% Medium & High &
% Limit the reference set to a representative subset; use explicit labeling guidelines; store rationale with labels; (if feasible) double-label a subset and resolve disagreements through adjudication; report uncertainty where evidence is ambiguous. \\

% Prototype scope or implementation effort exceeds available time &
% Medium & High &
% Timebox features and deliver an end-to-end minimal pipeline early (import model $\rightarrow$ parse stories $\rightarrow$ trace $\rightarrow$ coverage report);
% treat advanced strategies (e.g., LLM-based similarity) as optional; keep adapters modular to avoid rework when schemas change. \\

% Overfitting to one SUT or one exploration configuration &
% Medium & Medium &
% Evaluate on multiple runs/configurations; avoid hard-coded SUT-specific rules; document which choices are tool/SUT-specific; optionally validate on a second small SUT only if time allows. \\

% Low interpretability or poor reproducibility (heuristics/AI parameters) &
% Medium & Medium &
% Prefer deterministic baselines; version-control code, configurations, and input artifacts; fix random seeds where applicable; record parameters/thresholds in the repo and in the thesis; provide ablation-style comparison when more than one strategy is used. \\

% \bottomrule
% \end{tabular}
% \end{table}

% \section{Risk Analysis}
% \label{sec:risk-analysis}

% This section identifies key risks that may affect the design, implementation, and evaluation of the proposed method, and outlines mitigation measures. Likelihood and impact are assessed qualitatively as \emph{Low/Medium/High}.

% \begin{description}
%     \item[R1 --- Limited exploration coverage in GUI State Models:]
%     GUI State Models (D3) may not cover all relevant user story steps due to the stochastic nature of scriptless exploration, leading to many stories being labeled \emph{not covered} even when the SUT supports the behavior. \\
%     \textbf{Likelihood:} High. \\
%     \textbf{Impact:} High. \\
%     \textbf{Mitigation:} Run multiple explorations and/or configurations (D4) and aggregate evidence across models; clearly interpret \emph{not covered} as \emph{not covered in the available model evidence}; prioritize stories with observable UI effects.

%     \item[R2 --- Ambiguity and variability in user story wording:]
%     User stories (D2) may be underspecified, use domain-specific language, or describe outcomes not directly observable in the GUI State Model, reducing traceability. \\
%     \textbf{Likelihood:} Medium. \\
%     \textbf{Impact:} Medium. \\
%     \textbf{Mitigation:} Define a structured representation (A1) with explicit step/outcome fields; introduce normalization rules and a glossary for key terms; document assumptions and handle ambiguous stories via documented interpretation rules.

%     \item[R3 --- Incomplete or noisy model information:]
%     GUI State Models may contain missing labels, unstable widget identifiers, duplicated states, or transitions that do not correspond cleanly to story steps, causing incorrect trace links (A2). \\
%     \textbf{Likelihood:} Medium. \\
%     \textbf{Impact:} High. \\
%     \textbf{Mitigation:} Apply preprocessing and filtering (e.g., deduplication heuristics, label normalization); prefer robust features (e.g., action type + stable UI attributes); include failure-mode analysis in evaluation.

%     \item[R4 --- Difficulty constructing a reliable reference set:]
%     Manual creation of the reference set (D5) may be time-consuming and subject to disagreement, threatening evaluation validity. \\
%     \textbf{Likelihood:} Medium. \\
%     \textbf{Impact:} High. \\
%     \textbf{Mitigation:} Limit D5 to a representative subset; use clear labeling guidelines; record rationale for labels; perform adjudication for disagreements; report inter-annotator agreement where feasible.

%     \item[R5 --- Prototype implementation complexity and time constraints:]
%     Integrating parsing, model access (e.g., graph database), tracing, and reporting in the prototype (A4) may exceed available time. \\
%     \textbf{Likelihood:} Medium. \\
%     \textbf{Impact:} High. \\
%     \textbf{Mitigation:} Deliver a minimal end-to-end pipeline early; implement A1--A3 incrementally; timebox optional enhancements (e.g., advanced similarity/LLM support); keep interfaces modular so alternative tracing strategies can be swapped.

%     \item[R6 --- Overfitting to a single SUT or configuration:]
%     A method tuned to one SUT (D1) or one exploration configuration (D4) may not generalize. \\
%     \textbf{Likelihood:} Medium. \\
%     \textbf{Impact:} Medium. \\
%     \textbf{Mitigation:} Evaluate on multiple runs/configurations where possible; include robustness checks; make design choices explicit and discuss external validity limitations.

%     \item[R7 --- Threats to interpretability and reproducibility:]
%     If parts of the tracing rely on heuristic thresholds or AI-based similarity, results may be harder to reproduce or explain. \\
%     \textbf{Likelihood:} Medium. \\
%     \textbf{Impact:} Medium. \\
%     \textbf{Mitigation:} Prefer deterministic baselines; document parameters and thresholds; version-control code and configurations; report ablations (baseline vs enhanced strategy) when applicable.
% \end{description}
